\section{Ideal circuits with a finite attempt budget}\label{sec:finite}
The preceding stopping rule fixes the number of accepted readouts and
controls expected cost. It does not impose a deterministic deadline.
We now hold the ideal response exactly fixed, introduce no gate noise,
and instead stop after $T\geq1$ independent whole-filter attempts.
Every attempt reserves $c_\psi+4n+1$ query units, including readout on a
rejected attempt. Thus a hard budget $B$ permits
$T_n(B)=\lfloor B/(c_\psi+4n+1)\rfloor$ attempts without overspending.
This conservative reservation is common to all candidates.

Let $w(x)$ be the implemented success weight, $z=\sum_i p_iw(\lambda_i)$
and $\mu=\sum_i p_i\lambda_iw(\lambda_i)/z$.
For the Gram experiment $w=\kappa^2P^2$.
The three outcomes of an attempt have probabilities
$1-z,z(1-\mu),z\mu$: rejection, accepted zero and accepted one.
Consequently
\[
 M\sim\operatorname{Bin}(T,z),\qquad K\mid M=m\sim\operatorname{Bin}(m,\mu).
\]
If $M=0$, the algorithm reports no estimate. Otherwise it returns
$\widehat e=K/M$, without clipping to the prior interval.
We report abstention and conditional risk separately; successful trials
are not silently substituted for a fixed number of executions.

\begin{proposition}[Finite-attempt law]\label{prop:finite-law}
For $0<z\leq1$, set $q=1-z$, $C_T(z)=1-q^T$ and
\[
 H_T(z)=\sum_{m=1}^T\binom Tm z^m q^{T-m}/m,\qquad
 h_T(z)=H_T(z)/C_T(z).
\]
Then
\begin{equation}\label{eq:finite-mse}
 \Pr(M=0)=q^T,\qquad
 R_T(z,\mu,e):=\mathbb E[(\widehat e-e)^2\mid M>0]
 = (\mu-e)^2+\mu(1-\mu)h_T(z).
\end{equation}
The failure probability at tolerance $E$, counting abstention as failure,
is exactly
\begin{equation}\label{eq:finite-failure}
 q^T+\sum_{m=1}^T\binom Tm z^m q^{T-m}
 \Pr\{\lvert\operatorname{Bin}(m,\mu)/m-e\rvert>E\}.
\end{equation}
The function $h_T$ is nonincreasing on $(0,1]$, with
$h_T(1)=1/T$ and $\lim_{z\downarrow0}h_T(z)=1$.
\end{proposition}
\begin{proof}
Condition on $M=m>0$. The sample mean has mean $\mu$ and variance
$\mu(1-\mu)/m$. Averaging its squared error proves
\eqref{eq:finite-mse}; conditioning the error event proves
\eqref{eq:finite-failure}. The positive-binomial probabilities form an
exponential family in $\log(z/(1-z))$. Differentiating a finite sum gives
\[
 h_T'(z)=\frac{\operatorname{Cov}(1/M,M\mid M>0)}{z(1-z)}
 =\frac{1-h_T(z)\mathbb E[M\mid M>0]}{z(1-z)}\leq0,
\]
where Jensen's inequality gives the last step. The endpoint limits
follow from the finite sum, since at small $z$ a nonzero count is
asymptotically one.
\end{proof}

These binomial facts and inverse moments are classical
\cite{Znid,Drakakis}; we do not claim to discover shot noise or its
$1/\sqrt T$ regime. In particular, replacing $M$ by $Tz$ is not an exact
finite-sample analysis. When $Tz$ is small, conditioning on a nonempty
sample makes $h_T$ approach one, while abstention approaches one too.
As $T$ grows at fixed $z$, sampling error vanishes but the squared bias
$(\mu-e)^2$ remains. MSE and a fixed-confidence error guarantee are
different criteria; \eqref{eq:finite-failure} evaluates the latter.

\begin{corollary}[Uniform confidence at a hard deadline]\label{cor:deadline}
If $z\geq s>0$ and $0\leq\mu-e\leq b<E$ uniformly, counting abstention
as failure gives
\begin{equation}\label{eq:deadline-confidence}
 \Pr\{\text{no output or }|\widehat e-e|>E\}
 \leq\min\{1,\,2[1-s(1-e^{-2(E-b)^2})]^T\}.
\end{equation}
Thus a sufficient deterministic attempt count at confidence $1-\delta$ is
\[
 T\geq\left\lceil\frac{\log(2/\delta)}
 {-\log[1-s(1-e^{-2(E-b)^2})]}\right\rceil.
\]
\end{corollary}
\begin{proof}
For $M=m>0$, a deviation of size $E$ from $e$ requires one of size
at least $E-b$ from $\mu$. Hoeffding bounds its probability by
$2e^{-2m(E-b)^2}$. At $m=0$, the same expression is two and hence
upper bounds the failure indicator. Average over the binomial count
and use $z\geq s$ in its generating function.
\end{proof}
This is a standard concentration/PGF consequence, not a new concentration
inequality. It translates a supplied bias/success certificate to a
hard-deadline confidence statement, which the expected-cost stopping
analysis alone does not provide.

\begin{theorem}[Three-energy reduction for conditional risk]\label{thm:three}
Fix $T$, $a$, $u$, and a continuous weight $0\leq w\leq1$ that is
positive on $[0,u]$. The maximum of \eqref{eq:finite-mse} over all finite
compatible spectra and initial weights is attained by a measure
\begin{equation}\label{eq:three}
 a\delta_e+(1-a)\{\theta\delta_x+(1-\theta)\delta_y\},
 \quad 0\leq e\leq u,\quad e\leq x,y\leq1,\quad0\leq\theta\leq1.
\end{equation}
Repeated energies are allowed, so the actual target mass can exceed $a$.
\end{theorem}
\begin{proof}
For any input with minimum $e$ and target mass at least $a$, reserve mass
$a$ at $e$ and write the remainder as a probability measure $r$ on the
original finite support, including $e$ if necessary. Hold its resulting
conditional mean $\mu$ fixed. Among remainder weights minimize
$\int w\,dr$ subject to
\begin{equation}\label{eq:moment-lp}
 \int dr=1,\qquad
 \int (x-\mu)w(x)\,dr(x)
 =\frac{a}{1-a}(\mu-e)w(e),\qquad r\geq0.
\end{equation}
The original measure is feasible. This finite linear program has an
optimal basic solution with at most two positive remainder weights,
because it has at most two independent equality constraints. The new
input preserves $\mu$, decreases $z$, and therefore cannot decrease
$R_T$ by Proposition~\ref{prop:finite-law}. It has at most three energies.
Conversely every measure in \eqref{eq:three} is allowed. The parameter
set in \eqref{eq:three} is compact and $z\geq a\min_{[0,u]}w>0$, so its
continuous risk attains its maximum.
\end{proof}

This is a task-specific application of established moment-set and
basic-solution arguments \cite{Winkler,Pinelis}. It supplies an exact
low-dimensional characterization for this finite-risk objective, rather
than equating the separate worst bias, variance and acceptance witnesses.
It does not assert that three energies are always necessary, or that the
same reduction holds for the lattice-valued failure probability in
\eqref{eq:finite-failure}, a clipped estimator, or a different stopping rule.

\paragraph{A computable uniform selection bound.}
Suppose $z\geq s>0$ and $0\leq\mu-e\leq b$ uniformly. With
$v_b=t_b(1-t_b)$, $t_b=\min\{u+b,1/2\}$, one has
\begin{equation}\label{eq:finite-upper}
 R_T\leq b^2+v_bh_T(s),\qquad
 \Pr(M=0)\leq(1-s)^T.
\end{equation}
For endpoint-nulling polynomials the abstention bound is attained at the
minimum target transmission and the upper endpoint. The MSE bound need
not be attained, since its separate worst quantities can be incompatible.
To evaluate a rigorous upper bound without summing $T$ terms, put
\[
 A_T=\frac{1-(1-s)^{T+1}}{(T+1)s},\quad
 B_T=\frac{1-[1+(T+1)s](1-s)^{T+1}}{(T+1)(T+2)s^2}.
\]
Then
\begin{equation}\label{eq:inverse-upper}
 h_T(s)\leq\min\left\{1,
 \frac{A_T+3B_T-\frac52(1-s)^T}{1-(1-s)^T}\right\}.
\end{equation}
Indeed, $1/m\leq1/(m+1)+3/[(m+1)(m+2)]$ for $m\geq1$;
the full binomial expectations of these shifted inverse powers are
$A_T,B_T$, and removing $m=0$ subtracts $5(1-s)^T/2$.
All inputs in the implementation are rational. Directed 180-bit repeated
squaring encloses powers in \eqref{eq:inverse-upper}; numerical quadrature
is used only for auditing the exact law, never for certified selection.

For a finite candidate catalog, minimize this certified risk bound at
$T_n(B)$ subject to certified abstention at most $0.05$.
Include the deterministic prior midpoint, whose worst MSE is $u^2/4$
and whose abstention probability is zero. This gives a prior-only
finite-budget rule that can select no quantum work at small budgets.
It is optimal only for the computed bound within the supplied catalog,
not a new globally optimal quantum estimation algorithm.
The resource cautions of postselected metrology remain applicable
\cite{Combes}: conditional precision alone is not an advantage if the
discarded executions are omitted.
